\begin{figure}[h]
    \centering
    \begin{tikzpicture}[>=stealth, shorten >=1pt, auto, node distance=3cm, thick, main node/.style={circle, draw, font=\sffamily\Large\bfseries}]
    
      % Render the adjacency matrix
        \matrix[matrix of math nodes,
              left delimiter={[},
              right delimiter={]},
              %below = 1cm of A,
              nodes={minimum width=3em, minimum height=3em, outer sep=0pt},
              column sep=-\pgflinewidth,
              row sep=-\pgflinewidth] (M1) { 
                                            1 & 1 & 1 & 1 & 1 &\\
                                            };
        \node[left] at (M1.west) {Value List..};                                    
    \end{tikzpicture}

    \begin{tikzpicture}[>=stealth, shorten >=1pt, auto, node distance=3cm, thick, main node/.style={circle, draw, font=\sffamily\Large\bfseries}]

        \matrix[matrix of math nodes,
              left delimiter={[},
              right delimiter={]},
            %below = M1,
              nodes={minimum width=3em, minimum height=3em, outer sep=0pt},
              column sep=-\pgflinewidth,
              row sep=-\pgflinewidth] (M2) { 
                                            0 & 2 & 2 & 3 & 4 & 5\\
                                            };
        \node[left] at (M2.west) {Row Index.............}; 
    \end{tikzpicture}

    \begin{tikzpicture}[>=stealth, shorten >=1pt, auto, node distance=3cm, thick, main node/.style={circle, draw, font=\sffamily\Large\bfseries}]

        \matrix[matrix of math nodes,
              left delimiter={[},
              right delimiter={]},
              %below = M2,
              nodes={minimum width=3em, minimum height=3em, outer sep=0pt},
              column sep=-\pgflinewidth,
              row sep=-\pgflinewidth] (M3) { 
                                            1 & 2 & 2 & 2 & 0 & \\
                                            };
        \node[left] at (M3.west) {Col Index...};
    
    \end{tikzpicture}
    \caption{The \ac{CSR} format uses three lists to represent the graph in Figure~\ref{Figure:graph-example}.}
    \label{Figure:csr}
    \end{figure}